\scanpage{131}
\mspage{65}
locale autour de $x$ --- elle doit être essentiellement « triviale », que ma présomption s'avère vraie ou fausse --- par contre $1^\circ)$ est une question de nature globale sur $U$, et sans doute d'\deleted{\unsure{autre}} aucune façon triviale.

\deleted{Avant de \unsure{co|ce}} La validité de ces énoncés \deleted{signifierait} \inserted{impliquerait} donc, pour les sections de $\mathcal E_U^{\bar L}$ sur $\mathcal E_K^{\bar L}$ associées à un anneau de valuation de $L/K$, \inserted{de corps résiduel $K$}, que l'image de $\mathcal E_K^{\bar L}$ doit normaliser un \deleted{$I_x$} sous-groupe $\unsure{\Delta_x|\bar\Delta_x}\simeq T^{I(x)}$ de $\pi_{L/K}^{\bar L}$, qui est non trivial si la valuation n'a pas de centre sur $U$, i.e. si la section n'est pas associée à un pt $K$-rationnel de $U$, ce qui est justement dans le sens des conjectures (qui pouvaient d'abord sembler hâtives !) du § $2$.

Avant de me tourner vers l'étude des questions $1^\circ)$ et $2^\circ)$ (et \inserted{aussi} des questions de normalisation et de centralisation soulevées
\margin{130}
